\subsection{The randomized obligatory endpoint}
\label{sec:random-obligatory}

We now pass from the finite-cap randomized curve to the unbounded obligatory
endpoint.  The adversary fixes the labeled vector $p\in[0,\infty)^n$ before the algorithm
draws its random bits.  The multiset is not announced, and no upper bound on
the processing times is assumed.  The finite-cap zero--two lower bound
transfers directly.  The upper bound needs one additional uniformization:
a fixed cutoff and fallback protect the finite-cap stationary policy from an
arbitrarily large value missed by its sample.  This subsection proves
\cref{thm:rand-obligatory},
the rigorous randomized part of the informal obligatory-testing result
\cref{thm:intro-obligatory} from the introduction.

\begin{theorem}[Randomized obligatory testing]
\label{thm:rand-obligatory}
There is a universal randomized nonanticipating algorithm $\mathcal A_n$,
depending only on $n$, such that every fixed vector
$p\in[0,\infty)^n$ satisfies
\begin{equation}
 \EE\cost(\mathcal A_n,p)
 \le \frac43\cost(\OPT,p)+20378n^{7/4}
 \le\left(\frac43+40756n^{-1/4}\right)\cost(\OPT,p).
 \label{eq:rand-obligatory-upper}
\end{equation}
Conversely, for every randomized algorithm $\mathcal A$ and every $\eps>0$,
arbitrarily large even $n$ admit a fixed labeling of $n/2$ zero jobs and
$n/2$ jobs of processing time two such that
\begin{equation}
 \EE\cost(\mathcal A,p)
 \ge\left(\frac43-\eps\right)\cost(\OPT,p),
 \qquad
 \cost(\OPT,p)=\frac34n^2+n.
 \label{eq:rand-obligatory-lower}
\end{equation}
Hence the exact randomized size-asymptotic competitive ratio against an
oblivious adversary is $\RrandOT=4/3$.
\end{theorem}

For an announced multiset, the algorithmic core is the processing-time
prefix with maximum completions per unit discovery work.  Testing jobs in a
uniformly random order gives an exact pair identity and a $4/3$ algebraic
certificate.  The universal algorithm learns an approximately separated
prefix from a sublinear sample; a fixed cutoff and a fallback schedule
protect it from unbounded values missed by the sample.  The matching lower
bound is inherited from the zero--two endpoint of the finite-cap curve.  We
first prove the uniform upper bound and then give the short transfer.

\subsubsection{Normalization of the offline pair identity}

For a fixed vector $p$, use $D_p$ and its minimum kernel $K_p$ from
\cref{eq:stationary-kernel-split}.  The specialization in
\cref{eq:model-offline-coefficients} and the exact linear correction are
\begin{equation}
 \Omega_{\mathsf{OT}}(D_p)=\frac{1+K_p}{2},
 \qquad
 \Delta_{\mathsf{OT}}(p)=\frac{n+\sum_i p_i}{2}.
 \label{eq:rand-KO}
\end{equation}
The empirical effective-length
identity \cref{eq:empirical-offline-model}, with $\ell(p)=1+p$, gives
\begin{equation}
 \OPT(p)=n^2\Omega_{\mathsf{OT}}(D_p)+\Delta_{\mathsf{OT}}(p).
 \label{eq:rand-opt-normalized}
\end{equation}
The linear correction $\Delta_{\mathsf{OT}}(p)$ will be kept rather than
hidden: this makes the analysis valid for unbounded processing times.

\subsubsection{Stationary prefixes}

Use the processing-time-prefix notation
\cref{eq:stationary-prefix-notation,eq:stationary-kernel-split} and the
coefficient $\Psi_{D_p}(E)$ from
\cref{eq:stationary-leading-coefficient}.  We require that $E$ contain every
zero job; otherwise a supposedly deferred zero completes at its test.  The
stationary policy and its exact cost are supplied by the common pair identity
\cref{lem:stationary-prefix-cost}.

For a nonempty prefix define its completion density by
\begin{equation}
 \rho(E)=\frac{a}{1+m}.
 \label{eq:rand-density}
\end{equation}
Choose a maximum-density prefix and put
\begin{equation}
 \theta=\frac{1+m}{a}=\frac1{\rho(E)}.
 \label{eq:rand-theta}
\end{equation}
All zeros belong to every maximizer: adding a zero increases the numerator
without increasing the denominator.
The empirical prefix corollary \cref{cor:empirical-density-prefix} lets us
choose a maximizer with
\begin{equation}
 \theta=\tau_{D_p},\qquad
 p_i\le\theta\ (i\in E),\qquad p_i\ge\theta\ (i\notin E),
 \label{eq:rand-density-threshold}
\end{equation}
where tied occurrences may be assigned to either side.

The next identity is the algebraic source of the constant $4/3$.

\begin{lemma}[Exact and robust certificates]
\label{lem:rand-four-thirds-certificate}
For every ordered prefix split with $a>0$ and
$\theta=(1+m)/a$,
\begin{equation}
 2\theta\bigl(4\Omega_{\mathsf{OT}}(D_p)-3\Psi_{D_p}(E)\bigr)
  =(\theta-2)^2
   +4(\theta K_E-m^2)
   +\theta(K_T-\theta h^2).
 \label{eq:rand-certificate}
\end{equation}
Consequently:
\begin{enumerate}[label=\textup{(\roman*)}]
\item if early values are at most $\theta$ and late values are at least
      $\theta$, then $\Psi_{D_p}(E)\le4\Omega_{\mathsf{OT}}(D_p)/3$;
\item if early values are at most $\theta+s$ and late values are at least
      $\theta-s$, where $s\ge0$, then
      \begin{equation}
       \Psi_{D_p}(E)\le\frac43\Omega_{\mathsf{OT}}(D_p)+\frac23s.
       \label{eq:rand-robust-certificate}
      \end{equation}
\end{enumerate}
\end{lemma}

\begin{proof}
Expansion using $m=a\theta-1$, $h=1-a$, and
\cref{eq:stationary-kernel-split} gives \cref{eq:rand-certificate}.  Under exact
separation, $xy\le\theta\min\{x,y\}$ for early $x,y$, hence
$m^2\le\theta K_E$; for late $x,y$, one has
$\min\{x,y\}\ge\theta$, hence $K_T\ge\theta h^2$.  This proves (i).

Under the relaxed inequalities,
\[
 \theta K_E-m^2\ge-sK_E,
 \qquad K_T-\theta h^2\ge-sh^2.
\]
Moreover $K_E\le am\le a^2\theta$.  Substitution into
\cref{eq:rand-certificate} yields
\[
 \Psi_{D_p}(E)
 \le\frac43\Omega_{\mathsf{OT}}(D_p)+\frac23sa^2+\frac16sh^2
 \le\frac43\Omega_{\mathsf{OT}}(D_p)+\frac23s,
\]
because $(2/3)a^2+(1/6)(1-a)^2\le2/3$ on $[0,1]$.
\end{proof}

Combining the exact certificate with
\cref{lem:stationary-prefix-cost,eq:rand-opt-normalized} already gives a finite
statement useful below.

\begin{corollary}[Announced multiset]
\label{cor:rand-announced-four-thirds}
If the multiset is announced, the stationary policy for a maximum-density
prefix satisfies
\[
 \EE\ALG_{\Stat}(p,E)\le\frac43\OPT(p).
\]
\end{corollary}

Indeed, the correction in \cref{eq:stationary-prefix-cost} is at most
$\Delta_{\mathsf{OT}}(p)$.
Equality in the quadratic coefficient of \cref{eq:rand-certificate} forces
$\theta=2$, zero mass $1/2$, and all positive mass at two.  This predicts the
lower-bound instance before any probabilistic argument is used.

We also use the following tail-safe estimate.  Direct expansion gives
\begin{equation}
 \Psi_{D_p}(E)-\Omega_{\mathsf{OT}}(D_p)
 =\frac{h+am-K_E}{2}.
 \label{eq:rand-crude-identity}
\end{equation}
If all early values are at most $R\ge1$, then $m\le aR$ and therefore
\begin{equation}
 \Psi_{D_p}(E)\le \Omega_{\mathsf{OT}}(D_p)+\frac R2.
 \label{eq:rand-bounded-early}
\end{equation}

\subsubsection{Learning the prefix from a sublinear sample}

Fix the universal cutoff
\begin{equation}
 B=32.
 \label{eq:rand-cutoff}
\end{equation}
For an integer $d$, let $\eta=B/d$.  Keep zero as a separate category,
partition $(0,B]$ into the consecutive bins
$((b-1)\eta,b\eta]$, and add one overflow category for values above $B$.
A finite positive value is rounded upward to the endpoint $q_b=b\eta$.

Let $S$ be a uniformly random subset of $k$ labels.  If $\mu_b$ and
$\widehat\mu_b$ are the population and sample fractions of category $b$, put
\begin{equation}
 \Delta=\sum_b|\widehat\mu_b-\mu_b|.
 \label{eq:rand-hist-error}
\end{equation}

\begin{lemma}[Histogram error]
\label{lem:rand-histogram}
For $1\le k<n$,
\[
 \EE\Delta\le\sqrt{\frac{d+2}{k}}.
\]
\end{lemma}

\begin{proof}
This is \cref{eq:without-replacement-histogram}, with the zero bin and the
overflow bin included among the $d+2$ categories.
\end{proof}

For a finite category prefix $J$, define
\[
 \widehat a(J)=\sum_{b\in J}\widehat\mu_b,\qquad
 \widehat m(J)=\sum_{b\in J}q_b\widehat\mu_b,\qquad
 \widehat\rho(J)=\frac{\widehat a(J)}{1+\widehat m(J)}.
\]
Let $\widehat\rho$ be the maximum and
$\widehat\theta=1/\widehat\rho$.  The learned early set is not an arbitrary
maximizing category prefix.  It is its \emph{threshold closure}
\begin{equation}
 \widehat J=\{\text{finite categories }b:q_b\le\widehat\theta\}.
 \label{eq:rand-threshold-closure}
\end{equation}

\begin{lemma}[Threshold closure]
\label{lem:rand-threshold-closure}
The prefix \cref{eq:rand-threshold-closure} has the same maximum sample
density.  Every category in it has endpoint at most $\widehat\theta$, and
every finite category outside it has endpoint greater than
$\widehat\theta$, including categories with zero sample count.
\end{lemma}

\begin{proof}
The removal/addition calculation in
\cref{cor:empirical-density-prefix} applies to every positive-mass sample
category.  Empty categories do not change the sample statistics, and adding
mass at equality preserves the identity
$1+\widehat m=\widehat a\widehat\theta$.  Closing the prefix therefore
preserves optimality and makes the two pointwise assertions tautological.
\end{proof}

For large $n$, set
\begin{equation}
 k=\lfloor n^{3/4}\rfloor,\qquad
 d=\lfloor n^{1/4}\rfloor,\qquad \eta=B/d.
 \label{eq:rand-parameters}
\end{equation}
The algorithm draws a uniform permutation and first tests its first $k$ jobs,
leaving positive sample jobs pending.  It computes the closed sample prefix.
If
\begin{equation}
 \widehat\rho<\frac2B,
 \label{eq:rand-fallback-test}
\end{equation}
it enters \emph{fallback mode}: test every remaining job and then process all
positive jobs in SPT order.  Otherwise it enters \emph{learned mode}: process
positive sampled early jobs in SPT order; test the remaining jobs in their
random order, processing early outcomes immediately; and finish the deferred
tail in SPT order.  Overflow is never early.

\subsubsection{Good samples, fallback, and bad-event excess}

Put
\begin{equation}
 \gamma=\frac1{B(B+1)},\qquad G=\{\Delta\le\gamma\}.
 \label{eq:rand-good-event}
\end{equation}
Markov's inequality, recalled in \cref{sec:probabilistic-tools}, and
\cref{lem:rand-histogram} give
\begin{equation}
 \Prob(G^c)\le B(B+1)\sqrt{\frac{d+2}{k}}.
 \label{eq:rand-bad-probability}
\end{equation}

Assume $G$ and learned mode.  Let $E$ be the population set selected by the
closed category prefix, and let $a,m,\theta=(1+m)/a$ be its true statistics.
Since $\widehat\theta\le B/2$ and
$\widehat a=(1+\widehat m)/\widehat\theta\ge2/B$, we have $a\ge1/B$.  For
this fixed category set,
\[
 |a-\widehat a|\le\Delta,\qquad
 |m-\widehat m|\le B\Delta+\eta.
\]
Consequently
\begin{equation}
 |\theta-\widehat\theta|
 \le\frac32B^2\Delta+B\eta.
 \label{eq:rand-threshold-transfer}
\end{equation}
Every early actual value is at most $\widehat\theta$, while every late value
is greater than $\widehat\theta-\eta$.  Thus the robust certificate applies
with
\begin{equation}
 s=\frac32B^2\Delta+B\eta+\eta.
 \label{eq:rand-robust-s}
\end{equation}

The sample-first order is not exactly the fully stationary comparator, but
the difference is only lower order.  Conditional on the unordered sample,
write $e_S,e_R$ for the early counts and
\[
 W_S=k+\sum_{i\in E\cap S}p_i,\qquad
 W_R=n-k+\sum_{i\in E\setminus S}p_i.
\]
Comparing the two independent random internal orders with one fully random
order, the early-cost difference is
\[
 \frac{e_RW_S-e_SW_R}{2}
 \le\frac12nk(1+B/2).
\]
Postponing positive sampled early jobs until all sample tests finish adds at
most $k^2(1+B/2)$.  The late tail is identical because every zero is early.
Therefore, on a good learned sample,
\begin{equation}
 \EE[\ALG\mid S]
 \le\frac43\OPT
 +n^2\!\left[B^2\Delta+\frac23(B+1)\eta\right]
 +(1+B/2)(nk/2+k^2).
 \label{eq:rand-good-learned}
\end{equation}

On a good fallback sample, comparison of sample and population densities
shows that the true maximum-density reciprocal satisfies
$\tau_{D_p}\ge B/4=8$.
The exact threshold inequalities imply
\begin{equation}
 K_p\ge\frac{(\tau_{D_p}-1)^2}{\tau_{D_p}},
 \qquad
 \Omega_{\mathsf{OT}}(D_p)
 \ge\frac{\tau_{D_p}-1+1/\tau_{D_p}}{2}\ge\frac{57}{16}.
 \label{eq:rand-fallback-opt-large}
\end{equation}
For every test order, test-all-then-SPT satisfies
\begin{equation}
 \ALG_{\rm fb}\le\OPT+\frac{n(n-1)}2.
 \label{eq:rand-fallback-crude}
\end{equation}
Indeed, if there are $z$ zeros, their total displacement from the first $z$
SPT positions is at most $zn-z(z+1)/2$, while delaying the positive SPT tail
adds $(n-z)(n-z-1)/2$; the sum is $n(n-1)/2$.  Combining
\cref{eq:rand-fallback-opt-large,eq:rand-fallback-crude} gives
$\ALG_{\rm fb}<4\OPT/3$.

Because processing times are unbounded, a bad-event estimate must control
excess over $4\OPT/3$, not raw cost.  In learned mode every early value is at
most $B/2$, so \cref{eq:rand-bounded-early} bounds the stationary excess by
$(B/4)n^2$; in fallback mode \cref{eq:rand-fallback-crude} bounds it by
$n^2/2$.  Averaging \cref{eq:rand-good-learned} and these bad-event bounds,
and using \cref{eq:rand-bad-probability}, yields the explicit uniform estimate
\begin{equation}
 \EE\ALG
 \le\frac43\OPT
 +n^2\!\left[9472\sqrt{\frac{d+2}{k}}+\frac{704}{d}\right]
 +17(nk/2+k^2).
 \label{eq:rand-explicit-upper}
\end{equation}
For $n\ge12^4$, the floors in \cref{eq:rand-parameters} give
$\sqrt{(d+2)/k}\le2n^{-1/4}$, $1/d\le2n^{-1/4}$, and
$k\le n^{3/4}$.  The right-hand side of
\cref{eq:rand-explicit-upper} is at most
\begin{equation}
 \frac43\OPT+20378n^{7/4}.
 \label{eq:rand-final-upper}
\end{equation}
For smaller $n$, test-all-then-SPT is covered by the same deliberately loose
constant.  Finally $\OPT\ge n^2/2$ proves the multiplicative form in
\cref{eq:rand-obligatory-upper}.

\subsubsection{Transfer of the zero--two lower bound}

Fix the finite cap $u=25/4$.  Any randomized obligatory algorithm can be
viewed as a revealing-optimization algorithm at this cap that simply never
uses raw execution.  Apply \cref{lem:rcu-zero-two-lower}.  Its fixed hard
input has only values zero and two, so every revealing effective length is
\[
 \min\{25/4,1+p_j\}=1+p_j.
\]
The online transcript and the offline optimum are therefore exactly the
same in the finite-cap and obligatory models.  For even $n$, the common
offline value is
\[
 \OPT=\frac34n^2+n,
\]
whereas the selected labeling has expected online cost
$n^2-o(n^2)$.  Its ratio tends to $4/3$, which proves
\cref{eq:rand-obligatory-lower}.  Together with
\cref{eq:rand-final-upper}, this completes the proof of
\cref{thm:rand-obligatory}.
